\documentclass[10pt,a4paper]{amsart}
\usepackage{amsmath,amssymb,amsthm,graphicx}
\usepackage{verbatim}
\newenvironment{reference}{\comment}{\endcomment}
%\newenvironment{reference}{}{}
\newenvironment{slogan}{\comment}{\endcomment}
\newenvironment{history}{\comment}{\endcomment}

% For commutative diagrams we use Xy-pic
\usepackage[all]{xy}

% We use 2cell for 2-commutative diagrams.
\xyoption{2cell}
% Theorem environments.
%
\theoremstyle{plain}
\newtheorem{theorem}[subsection]{Theorem}
\newtheorem{proposition}[subsection]{Proposition}
\newtheorem{lemma}[subsection]{Lemma}

\theoremstyle{definition}
\newtheorem{definition}[subsection]{Definition}
\newtheorem{example}[subsection]{Example}
\newtheorem{exercise}[subsection]{Exercise}
\newtheorem{situation}[subsection]{Situation}

\theoremstyle{remark}
\newtheorem{remark}[subsection]{Remark}
\newtheorem{remarks}[subsection]{Remarks}

\numberwithin{equation}{subsection}

% Macros
%
\def\lim{\mathop{\mathrm{lim}}\nolimits}
\def\colim{\mathop{\mathrm{colim}}\nolimits}
\def\Spec{\mathop{\mathrm{Spec}}}
\def\Hom{\mathop{\mathrm{Hom}}\nolimits}
\def\Ext{\mathop{\mathrm{Ext}}\nolimits}
\def\SheafHom{\mathop{\mathcal{H}\!\mathit{om}}\nolimits}
\def\SheafExt{\mathop{\mathcal{E}\!\mathit{xt}}\nolimits}
\def\Sch{\mathit{Sch}}
\def\Mor{\mathop{\mathrm{Mor}}\nolimits}
\def\Ob{\mathop{\mathrm{Ob}}\nolimits}
\def\Sh{\mathop{\mathit{Sh}}\nolimits}
\def\NL{\mathop{N\!L}\nolimits}
\def\CH{\mathop{\mathrm{CH}}\nolimits}
\def\proetale{{pro\text{-}\acute{e}tale}}
\def\etale{{\acute{e}tale}}
\def\QCoh{\mathit{QCoh}}
\def\Ker{\mathop{\mathrm{Ker}}}
\def\Im{\mathop{\mathrm{Im}}}
\def\Coker{\mathop{\mathrm{Coker}}}
\def\Coim{\mathop{\mathrm{Coim}}}

% Boxtimes
%
\DeclareMathSymbol{\boxtimes}{\mathbin}{AMSa}{"02}

%
% Macros for moduli stacks/spaces
%
\def\QCohstack{\mathcal{QC}\!\mathit{oh}}
\def\Cohstack{\mathcal{C}\!\mathit{oh}}
\def\Spacesstack{\mathcal{S}\!\mathit{paces}}
\def\Quotfunctor{\mathrm{Quot}}
\def\Hilbfunctor{\mathrm{Hilb}}
\def\Curvesstack{\mathcal{C}\!\mathit{urves}}
\def\Polarizedstack{\mathcal{P}\!\mathit{olarized}}
\def\Complexesstack{\mathcal{C}\!\mathit{omplexes}}
% \Pic is the operator that assigns to X its picard group, usage \Pic(X)
% \Picardstack_{X/B} denotes the Picard stack of X over B
% \Picardfunctor_{X/B} denotes the Picard functor of X over B
\def\Pic{\mathop{\mathrm{Pic}}\nolimits}
\def\Picardstack{\mathcal{P}\!\mathit{ic}}
\def\Picardfunctor{\mathrm{Pic}}
\def\Deformationcategory{\mathcal{D}\!\mathit{ef}}
\usepackage{cleveref}
\crefname{originaltheorem}{original theorem}{original theorems}
\Crefname{originaltheorem}{Original theorem}{Original theorems}
\makeatletter
\newlabel{algebra-theorem-lazard}{{058G}{source}{}{sourcepointer-1}{}}
\newlabel{algebra-theorem-lazard@cref}{{[originaltheorem][0][]058G}{[1][0][]source}{}{}{}}
\makeatother
\crefname{originalreference}{original reference}{original references}
\Crefname{originalreference}{Original reference}{Original references}
\makeatletter
\newlabel{lemma-iso-go-up}{{083I}{source}{}{sourcepointer-2}{}}
\newlabel{lemma-iso-go-up@cref}{{[originalreference][0][]083I}{[1][0][]source}{}{}{}}
\makeatother
\crefname{originalreference}{original reference}{original references}
\Crefname{originalreference}{Original reference}{Original references}
\makeatletter
\newlabel{lemma-iso-limits}{{083J}{source}{}{sourcepointer-3}{}}
\newlabel{lemma-iso-limits@cref}{{[originalreference][0][]083J}{[1][0][]source}{}{}{}}
\makeatother
\crefname{originalreference}{original reference}{original references}
\Crefname{originalreference}{Original reference}{Original references}
\makeatletter
\newlabel{situation-somewhat-closed}{{0CWA}{source}{}{sourcepointer-4}{}}
\newlabel{situation-somewhat-closed@cref}{{[originalreference][0][]0CWA}{[1][0][]source}{}{}{}}
\makeatother
\crefname{originallemma}{original lemma}{original lemmas}
\Crefname{originallemma}{Original lemma}{Original lemmas}
\makeatletter
\newlabel{spaces-properties-lemma-cover-by-union-affines}{{03FX}{source}{}{sourcepointer-5}{}}
\newlabel{spaces-properties-lemma-cover-by-union-affines@cref}{{[originallemma][0][]03FX}{[1][0][]source}{}{}{}}
\makeatother
\title{The base change where a map becomes zero is cut out by a closed subscheme, of finite presentation in the finite type case}
\author{The Stacks Project Authors}
\date{stacks-project:spaces-flat.tex:local-snapshot-20261012; GNU FDL 1.2 or later}
\begin{document}
\maketitle
\noindent\textit{Editorial note.} All delivered mathematical text is the author's own, copied from the pinned Stacks Project chapter file spaces-flat.tex; no statement or proof step was rewritten, reordered or removed.\par\medskip
\noindent\textit{Editorial note.} The counted proof is the author's complete argument: a reduction to affine pieces, a value by value construction of the killing ideal through Lazard's theorem, and a limit argument producing a finitely generated ideal in the finite type case\par\medskip
\noindent\textit{Editorial note.} The label lemma-F-zero-somewhat-closed names the counted statement, and the proof environment immediately following it in the source is delivered verbatim as its proof.\par\medskip
\noindent\textit{Editorial note.} The counted statement is carried by the author's own lemma environment in the chapter source.\par\medskip
\noindent\textit{Allowed standard prerequisites (editorial).}\begin{itemize}
\item Quasi-coherent modules on algebraic spaces, covers by disjoint unions of affine schemes, and the go-up lemma allowing the statement to be checked after a surjective etale base change
\item Lazard's theorem that a flat module is a filtered colimit of free modules, and commutation of tensor product with colimits
\item The limit criterion that a vanishing statement stable under base change holds already at some stage of a filtered colimit of closed subschemes
\end{itemize}\medskip
\ifdefined\hypertarget\hypertarget{sourcepointer-1}{}\fi
\noindent\textit{Source reference.} Original theorem 058G: sources/stacks-chow/tags.txt, source lines 6779--6779. Prerequisite pointer invoked at proof line(s) 1604; it names a standard fact of a sibling chapter that the counted argument uses, and every step of the counted proof built on it is retained verbatim in the delivered body.\par\medskip
\ifdefined\hypertarget\hypertarget{sourcepointer-2}{}\fi
\noindent\textit{Source reference.} Original reference 083I: sources/stacks-spaces-flat/pinned-original-source.tex, source lines 1291--1291. Prerequisite pointer invoked at proof line(s) 1582; it names a standard fact of the same chapter that the counted argument uses, and every step of the counted proof built on it is retained verbatim in the delivered body.\par\medskip
\ifdefined\hypertarget\hypertarget{sourcepointer-3}{}\fi
\noindent\textit{Source reference.} Original reference 083J: sources/stacks-spaces-flat/pinned-original-source.tex, source lines 1335--1335. Prerequisite pointer invoked at proof line(s) 1622; it names a standard fact of the same chapter that the counted argument uses, and every step of the counted proof built on it is retained verbatim in the delivered body.\par\medskip
\ifdefined\hypertarget\hypertarget{sourcepointer-4}{}\fi
\noindent\textit{Source reference.} Original reference 0CWA: sources/stacks-spaces-flat/pinned-original-source.tex, source lines 1560--1560. Prerequisite pointer invoked at proof line(s) 1568; it names a standard fact of the same chapter that the counted argument uses, and every step of the counted proof built on it is retained verbatim in the delivered body.\par\medskip
\ifdefined\hypertarget\hypertarget{sourcepointer-5}{}\fi
\noindent\textit{Source reference.} Original lemma 03FX: sources/stacks-chow/tags.txt, source lines 4508--4508. Prerequisite pointer invoked at proof line(s) 1581; it names a standard fact of a sibling chapter that the counted argument uses, and every step of the counted proof built on it is retained verbatim in the delivered body.\par\medskip
\begin{lemma}
\label{lemma-F-zero-somewhat-closed}
In Situation \ref{situation-somewhat-closed}. Let $T \to S$
be a quasi-compact morphism of schemes such that the base change $u_T$ is
zero. Then exists a closed subscheme $Z \subset S$ such that
(a) $T \to S$ factors through $Z$ and (b) the base change $u_Z$ is zero.
If $\mathcal{F}$ is a finite type $\mathcal{O}_X$-module and
the scheme theoretic support of $\mathcal{F}$ is quasi-compact,
then we can take $Z \to S$ of finite presentation.
\end{lemma}

\begin{proof}
Let $U \to X$ be a surjective \'etale morphism of algebraic spaces
where $U = \coprod U_i$ is a disjoint union of affine schemes (see
Properties of Spaces, Lemma
\ref{spaces-properties-lemma-cover-by-union-affines}).
By Lemma \ref{lemma-iso-go-up} we see that we may
replace $X$ by $U$. In other words, we may assume that $X = \coprod X_i$
is a disjoint union of affine schemes $X_i$. Suppose that we can prove
the lemma for $u_i = u|_{X_i}$. Then we find a closed subscheme
$Z_i \subset S$ such that $T \to S$ factors through $Z_i$ and
$u_{i, Z_i}$ is zero. If
$Z_i = \Spec(R/I_i) \subset \Spec(R) = S$, then taking
$Z = \Spec(R/\sum I_i)$ works. Thus we may assume that
$X = \Spec(A)$ is affine.

\medskip\noindent
Choose a finite affine open covering $T = T_1 \cup \ldots \cup T_m$.
It is clear that we may replace $T$ by $\coprod_{j = 1, \ldots, m} T_j$.
Hence we may assume $T$ is affine. Say $T = \Spec(R')$.
Let $u : M \to N$ be the homomorphisms of $A$-modules
corresponding to $u : \mathcal{F} \to \mathcal{G}$.
Then $N$ is a flat $R$-module as $\mathcal{G}$ is flat over $S$.
The assumption of the lemma means that the composition
$$
M \otimes_R R' \to N \otimes_R R'
$$
is zero. Let $z \in M$. By Lazard's theorem
(Algebra, Theorem \ref{algebra-theorem-lazard}) and the fact
that $\otimes$ commutes with colimits we can find free $R$-module
$F_z$, an element $\tilde z \in F_z$, and a map $F_z \to N$ such that
$u(z)$ is the image of $\tilde z$ and $\tilde z$ maps to zero in
$F_z \otimes_R R'$. Choose a basis $\{e_{z, \alpha}\}$ of $F_z$ and write
$\tilde z = \sum f_{z, \alpha} e_{z, \alpha}$ with $f_{z, \alpha} \in R$.
Let $I \subset R$ be the ideal generated by the elements $f_{z, \alpha}$
with $z$ ranging over all elements of $M$.
By construction $I$ maps to zero in $R'$ and the elements $\tilde z$
map to zero in $F_z/IF_z$ whence in $N/IN$. Thus $Z = \Spec(R/I)$
is a solution to the problem in this case.

\medskip\noindent
Assume $\mathcal{F}$ is of finite type with quasi-compact scheme
theoretic support. Write $Z = \Spec(R/I)$.
Write $I = \bigcup I_\lambda$ as a filtered union of finitely generated
ideals. Set $Z_\lambda = \Spec(R/I_\lambda)$, so $Z = \colim Z_\lambda$.
Since $u_Z$ is zero, we see that $u_{Z_\lambda}$ is zero
for some $\lambda$ by Lemma \ref{lemma-iso-limits}.
This finishes the proof of the lemma.
\end{proof}

\end{document}
